% ---------------------------------------------------------------------------
% Introduction draft for the DIF-Gaussian extension paper.
% Base work: Lin et al., "Learning 3D Gaussians for Extremely Sparse-View
% Cone-Beam CT Reconstruction", MICCAI 2024 (arXiv:2407.01090).
% Rename \method{} in one place if a different name is chosen.
% Citation keys below are placeholders; verify all entries against the
% original papers before submission (suggested refs listed in the comment
% block at the end of this file).
% ---------------------------------------------------------------------------

\newcommand{\method}{DIF-Gaussian++}

\section{Introduction}
\label{sec:intro}

Cone-beam computed tomography (CBCT) plays a central role in image-guided
surgery, radiotherapy, and dental imaging, and is increasingly explored for
low-dose lung screening. In all these applications, radiation dose, scanning
time, and hardware cost scale with the number of acquired projections, which
motivates \emph{extremely sparse-view} protocols that reconstruct a
diagnostic-quality volume from only a handful of projections. However, the
sampling density for which filtered backprojection (FDK)
\cite{feldkamp1984practical} remains accurate is far below what such protocols
demand: under extremely sparse sampling, view-starvation artifacts---streaks,
angular blurring, and eroded boundaries---dominate the reconstructed volume.

Existing approaches fall into three broad families. \emph{Analytic and
model-based iterative methods} formulate reconstruction as regularized
inversion of the imaging model; compressed-sensing approaches with
hand-crafted priors such as total variation markedly reduce FDK artifacts
\cite{sidky2008image}, but require minutes to hours of per-patient
optimization and are sensitive to system-model mismatch and regularization
hyperparameters. \emph{Supervised deep learning methods} learn the mapping
from measurements to images feed-forward: most operate on individual slices
or projections, so residual artifacts remain inconsistent between slices and
3D streaks survive \cite{jin2017deep,kang2017deep}; learned iterative and 3D
variants partially mitigate this but inherit heavy inference cost or need
dense-view supervision \cite{adler2018learned}. \emph{Implicit neural
representations} optimize a continuous volumetric network per patient---either
as INR-based attenuation fields \cite{zang2021naf,cai2023neat,cai2024semiinr}
or, more recently, as radiative 3D Gaussians
\cite{zhang2024r2gaussian,cai2024xgaussian}---achieving excellent view
consistency, but fitting every new patient takes tens of minutes and none of
these representations generalizes beyond the case being optimized. A recent
line of work therefore targets \emph{generalizable} implicit-field

Our contributions are fourfold:
\begin{itemize}
  \item We identify cross-view evidence fusion---rather than the implicit
  field representation itself---as the bottleneck of generalizable Gaussian
  implicit-field reconstruction in the extremely sparse-view regime, and
  analyze why view-agnostic max fusion fails: it is geometry-blind,
  evidence-discarding, and spatially isolated.
  \item We propose a geometry-aware multiview transformer that lifts
  multi-scale 2D features onto shared 3D anchor tokens: per-anchor
  cross-attention over ray-encoded, FOV-masked observations from all views and
  scales, followed by windowed 3D self-attention with physical
  relative-position biases.
  \item We integrate this encoder with the Gaussian implicit field so that
  anisotropic Gaussian attributes and content features are predicted from
  geometry-fused tokens, preserving continuous, memory-bounded query decoding
  and generalization across patients and view counts.
  \item On the LUNA16-based C2RV-CBCT benchmark with only six projections,
  our method consistently outperforms the max-fusion baseline and a late
  geometry-attention variant, and systematic ablations isolate the
  contribution of each design choice (fusion mechanism, multi-scale
  observations, spatial window, token initialization and residual design).
  % TODO: verify this sentence against the final 400-epoch E400 tables before
  % submission; soften if some ablations do not show consistent gains.
\end{itemize}

% ---------------------------------------------------------------------------
% Suggested references (VERIFY bibliographic details before use):
%
%  feldkamp1984practical - L.A. Feldkamp, L.C. Davis, J.W. Kress,
%      "Practical cone-beam algorithm", J. Opt. Soc. Am. A, 1(6):612-619, 1984.
%  sidky2008image        - E.Y. Sidky, X. Pan, "Image reconstruction in
%      circular cone-beam computed tomography by constrained, total-variation
%      minimization", Phys. Med. Biol., 53(17):4777-4807, 2008.
%  jin2017deep           - K.H. Jin, M.T. McCann, E. Froustey, M. Unser,
%      "Deep convolutional neural network for inverse problems in imaging",
%      IEEE TIP, 26(9):4509-4522, 2017.
%  kang2017deep          - E. Kang et al., "Wavelet domain residual network
%      for low-dose X-ray CT reconstruction", arXiv:1703.01383, 2017.
%  adler2018learned      - J. Adler, O. Oktem, "Learned primal-dual
%      reconstruction", IEEE TMI, 37(6):1322-1332, 2018.
%  zang2021naf           - G. Zang et al., "Neural attenuation fields for
%      sparse-view CBCT reconstruction", Medical Image Analysis, 2023.
%  cai2023neat           - R. Cai et al., "NeAT: Neural adaptive tomography",
%      ICCV 2023.
%  cai2024semiinr        - R. Cai et al., "Structured uncertainty ... (Semi-INR)",
%      ICLR 2024.
%  zhang2024r2gaussian   - S. Zhang et al., "R2-Gaussian: Rectifying radiative
%      Gaussian splatting for tomographic reconstruction", NeurIPS 2024.
%  cai2024xgaussian      - R. Cai et al., "Radiative Gaussian splatting for
%      efficient X-ray novel view synthesis (X-Gaussian)", ECCV 2024.
%  lin2024dif            - Y. Lin, H. Wang, J. Chen, X. Li, "Learning 3D
%      Gaussians for extremely sparse-view cone-beam CT reconstruction",
%      MICCAI 2024 (arXiv:2407.01090).
%
% Optional additions: iterative reconstruction reviews (e.g., Nuyts et al.),
% TV/CS variants, DDO/dual-domain networks, Transformer-based recon
% (e.g., DuDoTrans), learned LIRE, CORAL, FBI.
% ---------------------------------------------------------------------------
